\documentclass[../preamble]{subfiles}
\begin{document}

\section{Functors}
\label{sec:functor}
\concept{functor}{Functor}
\depends{category}
\implements{haskell}{Cat.Functor}
\implements{agda}{Cat.Functor}
\implements{lisp}{functor}

Let $\Category{C}$ and $\Category{D}$ be categories.
A \newterm{functor} $\newmath{\Functor{F} \colon \Category{C} \lto \Category{D}}$
consists of the following data
(\stacksref{001B}; \nlabref{functor}):
\begin{enumerate}
  \item A map on objects
    \begin{equation}
      \begin{aligned}
        \Functor{F} \colon \Ob(\Category{C}) &\lto \Ob(\Category{D}) \\
        X &\lmapsto \Functor{F}(X).
      \end{aligned}
      \label{eq:functor-ob}
    \end{equation}
  \item For each pair of objects $X, Y \in \Ob(\Category{C})$,
    a map on hom-sets
    \begin{equation}
      \begin{aligned}
        \Functor{F}_{X,Y} \colon
          \Hom_{\Category{C}}(X, Y)
          &\lto
          \Hom_{\Category{D}}(\Functor{F}(X), \Functor{F}(Y)) \\
        f &\lmapsto \Functor{F}(f).
      \end{aligned}
      \label{eq:functor-hom}
    \end{equation}
\end{enumerate}

subject to the following axioms:
\begin{itemize}
  \item \textbf{Preservation of identities.}
    \axiom{functor}{preserve-identity}
    For each object $X \in \Ob(\Category{C})$,
    \begin{equation}
      \label{eq:functor-id}
      \Functor{F}(\id_X) = \id_{\Functor{F}(X)}.
    \end{equation}

  \item \textbf{Preservation of composition.}
    \axiom{functor}{preserve-composition}
    For all morphisms $f \colon X \to Y$ and $g \colon Y \to Z$ in $\Category{C}$,
    \begin{equation}
      \label{eq:functor-comp}
      \Functor{F}(g \circ f) = \Functor{F}(g) \circ \Functor{F}(f).
    \end{equation}
\end{itemize}

\begin{definition}[Identity functor]
  \label{def:identity-functor}
  \concept{identity-functor}{Identity Functor}
  \depends{functor}
  The \newterm{identity functor} on $\Category{C}$ is
  $\newmath{\Functor{Id}_{\Category{C}}} \colon \Category{C} \lto \Category{C}$,
  defined by
  \begin{equation}
    \begin{aligned}
      \Functor{Id}_{\Category{C}} \colon \Ob(\Category{C}) &\lto \Ob(\Category{C}) \\
      X &\lmapsto X,
    \end{aligned}
    \label{eq:id-functor-ob}
  \end{equation}
  and $\Functor{Id}_{\Category{C}}(f) = f$ on morphisms.
  It satisfies both functor axioms trivially.
\end{definition}

\begin{definition}[Composition of functors]
  \label{def:functor-composition}
  \concept{functor-composition}{Functor Composition}
  \depends{functor}
  Given functors $\Functor{F} \colon \Category{C} \lto \Category{D}$
  and $\Functor{G} \colon \Category{D} \lto \Category{E}$,
  their \newterm{composite}
  $\newmath{\Functor{G} \circ \Functor{F}} \colon \Category{C} \lto \Category{E}$
  is defined by
  \begin{align}
    \label{eq:functor-comp-ob}
    (\Functor{G} \circ \Functor{F})(X) &= \Functor{G}(\Functor{F}(X))
      \quad \text{on objects,} \\
    \label{eq:functor-comp-mor}
    (\Functor{G} \circ \Functor{F})(f) &= \Functor{G}(\Functor{F}(f))
      \quad \text{on morphisms.}
  \end{align}
\end{definition}

\begin{proposition}
  \label{prop:functor-comp-is-functor}
  The composite $\Functor{G} \circ \Functor{F}$ as defined above is a functor.
\end{proposition}

\begin{proof}
  Preservation of identities:
  \[
    (\Functor{G} \circ \Functor{F})(\id_X)
    = \Functor{G}(\Functor{F}(\id_X))
    = \Functor{G}(\id_{\Functor{F}(X)})
    = \id_{\Functor{G}(\Functor{F}(X))}
    = \id_{(\Functor{G} \circ \Functor{F})(X)}.
  \]
  Preservation of composition: for $f \colon X \to Y$ and $g \colon Y \to Z$,
  \[
    (\Functor{G} \circ \Functor{F})(g \circ f)
    = \Functor{G}(\Functor{F}(g \circ f))
    = \Functor{G}(\Functor{F}(g) \circ \Functor{F}(f))
    = \Functor{G}(\Functor{F}(g)) \circ \Functor{G}(\Functor{F}(f))
    = (\Functor{G} \circ \Functor{F})(g) \circ (\Functor{G} \circ \Functor{F})(f). \qedhere
  \]
\end{proof}

\begin{remark}
  Functor composition is associative,
  and $\Functor{Id}_{\Category{C}}$ is a two-sided unit.
  Consequently, small categories and functors form a category
  $\newmath{\Category{Cat}}$.
\end{remark}

\end{document}
